% ============================================================
% §5.3 Q3结果分析 — 服务定价与政府补贴 MILP (v2: F/H/J配置)
% 竞赛: 电工杯 B题 (2026)
% ============================================================

\section{Q3: 服务定价与政府补贴优化结果}

\subsection{最优定价方案}

基于§3.2构造的S3价格满意度Big-M线性化模型，在固化Q2选址方案F大型+H中型+J中型且10/10全覆盖后，以词典序目标求解MILP——优先最大化加权$S_3$满意度，次优先以弱激励$\varepsilon=10^{-4}$最大化营收。CBC在1秒内收敛至全局最优，全部6项服务均实现$S_3=1.000$平价区间，仅上门护理一项需从基准价30元降至12.64元，降幅57.9\%，以满足F站的8\%利润率紧约束。最优定价方案见表\ref{tab:q3_pricing}，定价对比与S3满意度见图\ref{fig:q3_pricing}。

\begin{figure}[H]
\centering
\includegraphics[width=0.92\textwidth]{figures/figure_q3_pricing_comparison.png}
\caption{基准价与最优定价对比。上门护理从30元降至12.64元是唯一的降价服务，降幅57.9\%，其余5项服务维持基准价。全部服务S3=1.000平价区间。}
\label{fig:q3_pricing}
\end{figure}

\begin{table}[htbp]
\centering
\caption{Q3最优定价方案}
\label{tab:q3_pricing}
\begin{tabular}{lccccc}
\toprule
\textbf{服务项目} & \textbf{基准价(元)} & \textbf{最优定价(元)} & \textbf{溢价率(\%)} & \textbf{S3} & \textbf{价格区间} \\
\midrule
助餐 & 10 & 10.00 & 0.0 & 1.000 & 平价 \\
日间照料 & 20 & 20.00 & 0.0 & 1.000 & 平价 \\
上门护理 & 30 & \textbf{12.64} & \textbf{-57.9} & 1.000 & 平价 \\
康复理疗 & 28 & 28.00 & 0.0 & 1.000 & 平价 \\
助浴 & 25 & 25.00 & 0.0 & 1.000 & 平价 \\
紧急救助 & 0 & 0.00 & 公益免费 & 1.000 & 公益免费 \\
\bottomrule
\end{tabular}
\end{table}

\subsection{站点利润率约束的绑定诊断}

表\ref{tab:q3_stations}给出了最优定价下的站点级运营指标。F站的利润率恰好触达8.0\%的监管红线，这是全系统唯一的紧约束。H站与J站的利润率仅为7.3\%，说明在S3最大化的词典序目标下两站尚存0.7个百分点的定价空间，但因$S_3$已全域触及1.000平价段上线，任何提价行为都会将对应服务的$S_3$从1.000降至0.900微溢价段，造成加权满意度目标函数的净损失。这一结果完美诠释了词典序目标的设计意图：\textbf{在满足所有站点可持续运营、利润率不超8\%的前提下，S3优先于营收}。

\begin{table}[htbp]
\centering
\caption{最优定价下的站点运营指标}
\label{tab:q3_stations}
\small
\begin{tabular}{lcccccc}
\toprule
\textbf{站点} & \textbf{规模} & \textbf{年营收} & \textbf{年补贴} & \textbf{年总成本} & \textbf{年利润} & \textbf{利润率} \\
 &  & \textbf{(万元)} & \textbf{(万元)} & \textbf{(万元)} & \textbf{(万元)} & \textbf{(\%)} \\
\midrule
F & 大型 & 1,967.1 & 93.6 & 1,908.1 & 152.6 & \textbf{8.0} \\
H & 中型 & 1,265.6 & 64.8 & 1,239.6 & 90.8 & 7.3 \\
J & 中型 & 1,226.3 & 64.8 & 1,203.2 & 87.9 & 7.3 \\
\midrule
\textbf{合计} & -- & 4,459.0 & 223.2 & 4,350.9 & 331.3 & 7.6 \\
\bottomrule
\end{tabular}
\end{table}

\subsection{三层交叉补贴机制分析}

\subsubsection{第一层：站内营利养公益}

Q3的全系统紧急救助月需求为4,940次，其中F站2,202次、H站1,399次、J站1,339次，单价0元公益免费，单次直接成本8元。全系统紧急救助年净支出为47.42万元——这笔纯公益支出由各站营利性服务的年总利润331.3万元完全吸收，交叉补贴率仅14.3\%。全覆盖带来的规模经济效应显著摊薄了公益服务的单位吸收成本。

\subsubsection{第二层：站间补贴帽的效率错配}

政府补贴规则为非紧急救助服务每单补贴2元，受站点日补贴上限约束。在三站满负荷运营的条件下，各站的理论日补贴额与实际上限的对比见表\ref{tab:q3_subsidy}。三个站点的有效补贴率均远低于理论的2元/次——F站0.73元/次、H站0.79元/次、J站0.81元/次。F站因规模最大、需求最多，63.5\%的潜在补贴被上限截断，效率损失最为严重，形成一种逆向再分配：大型站因服务人次多而触发补贴帽更早，有效补贴率反而最低。

\begin{table}[htbp]
\centering
\caption{各站点补贴池约束分析}
\label{tab:q3_subsidy}
\small
\begin{tabular}{lccccc}
\toprule
\textbf{站点} & \textbf{日非紧急需求} & \textbf{理论补贴} & \textbf{日补贴上限} & \textbf{实际补贴} & \textbf{有效补贴率} \\
 & \textbf{(次)} & \textbf{(元/日)} & \textbf{(元)} & \textbf{(元/日)} & \textbf{(元/次)} \\
\midrule
F & 3,559 & 7,118 & 2,600 & 2,600 & 0.73 \\
H & 2,291 & 4,582 & 1,800 & 1,800 & 0.79 \\
J & 2,221 & 4,442 & 1,800 & 1,800 & 0.81 \\
\bottomrule
\end{tabular}
\end{table}

\subsubsection{第三层：上门护理的"拉格朗日影子价格"效应}

在Q3的最优解中，上门护理从30元降至12.64元，降幅57.9\%，成为唯一的降价服务，而助餐虽频次最高、需求弹性最大却维持基准价不动。这一结果与站点规模结构直接相关：F站作为利润率紧约束的绑定者，其上门护理在所有服务中基准价最高，单位降价的利润率释放效应最大。从拉格朗日乘子的视角，上门护理的影子价格——降低1元对利润率松绑的边际贡献——在所有服务中最高，因此成为词典序S3最大化目标下的最优调节变量。

从Ramsey-Boiteux最优定价的视角看：在总利润率约束下，价格需求弹性最小的服务应承担最少的降价。上门护理的基准价30元在所有服务中最高，其过剩的利润空间使其成为$S_3$最大化目标下最优先被削减的对象。这一结果可理解为反向价格歧视——高频低价服务助餐每月118,443次受保护，低频高价服务上门护理每月19,608次承担调节压力，在社会福利函数中体现了利用最不敏感价格维度吸收监管约束的Ramsey原理\cite{ramsey1927}。补贴机制中"补需方"与"补供方"的效率权衡以及预算约束下的补贴策略优化，为本文三层交叉补贴的额度标定提供了制度经济学参照。

\subsubsection{交叉补贴机制的算法透明化}

三层交叉补贴的决策逻辑——形式化算法见附录\ref{alg:cross_subsidy}——体现为三阶段不可逆的递进流程：微利约束$\to$补贴效率$\to$定价弹性。

\subsection{加权综合满意度解构}

表\ref{tab:q3_satisfaction}将各小区的三维满意度分解为$S_1$（距离，已由Q2分配固定）、$S_2$（响应满意度，全部0.600因三站满负荷）、$S_3$（价格满意度，全部1.000因Q3最优定价维持平价）。加权综合满意度$S$的跨小区变异系数仅为0.035，区间[0.800, 0.880]，反映了全覆盖方案下满意度的空间均等化特征。

\begin{table}[htbp]
\centering
\caption{各小区加权综合满意度分解（Q2+Q3联合）}
\label{tab:q3_satisfaction}
\begin{tabular}{lcccccc}
\toprule
\textbf{小区} & \textbf{服务站} & $S_1$(距离) & $S_2$(响应) & $S_3$(价格) & $S$(综合) & \textbf{说明} \\
\midrule
A & J & 0.900 & 0.600 & 1.000 & 0.860 & 中距离，满意均好 \\
B & H & 0.900 & 0.600 & 1.000 & 0.860 & 中距离，满意均好 \\
C & F & 0.600 & 0.600 & 1.000 & 0.800 & 远距离S1触底 \\
D & J & 0.900 & 0.600 & 1.000 & 0.860 & 中距离，满意均好 \\
E & H & 0.600 & 0.600 & 1.000 & 0.800 & 边界距离S1触底 \\
F & F & 1.000 & 0.600 & 1.000 & 0.880 & 本小区自服务 \\
G & F & 0.750 & 0.600 & 1.000 & 0.830 & 中距离，满意均好 \\
H & H & 1.000 & 0.600 & 1.000 & 0.880 & 本小区自服务 \\
I & F & 0.600 & 0.600 & 1.000 & 0.800 & 远距离S1触底 \\
J & J & 1.000 & 0.600 & 1.000 & 0.880 & 本小区自服务 \\
\midrule
\textbf{平均} & -- & 0.825 & 0.600 & 1.000 & 0.845 & -- \\
\bottomrule
\end{tabular}
\end{table}

\subsection{服务可及性的三维透视}

\textbf{经济可及性。} Q1.3预诊断表明失能全10小区触发削减，缺口23.9\%--49.1\%。经Q3优化，上门护理从30元降至12.64元，降幅57.9\%，单次消费占比从30\%+压缩至12.7\%，大批因预算缺口被拦在门外的失能群体重新进入有效需求。每1元财政补贴撬动约2.4元降价空间。

\textbf{地理空间可及性。} 1000m服务半径对三类老人呈梯度分化：自理老人自主出行，1000m为软约束；半失能老人步行能力退化，边界小区存在可及性脆弱性；失能卧床老人100\%依赖上门送达，1000m约束转化为服务人员通勤成本。核心原则：\textbf{站点选址应向失能密度高的区域倾斜，而非简单追求几何中心性}。

\textbf{数字信息可及性。} 自理老人可通过智能终端主动触达，半失能依赖简化交互与子女代理，失能独居面临数字鸿沟。紧急救助作为唯一"被动触发"服务——一键呼叫——其公益免费属性确保数字弱势群体不失最后防线，全系统年净支出47.42万元可视为政府购买"数字公平"公共品的财政对价。
